\usepackage{fullpage}

\usepackage[dvipsnames,svgnames,table]{xcolor}
\usepackage[colorlinks=true,linkcolor=RawSienna,citecolor=ForestGreen]{hyperref}

\usepackage[utf8]{inputenc}

\usepackage{amsthm}
\usepackage{thmtools}
\usepackage[normalem]{ulem}
\usepackage{mathtools,bm}
\usepackage[inline]{enumitem}
\usepackage{amsmath,amsfonts,amssymb,amsthm,hyperref,algorithm,eqparbox,array,ragged2e,stmaryrd,bbm}
\usepackage{xspace}
\usepackage{xcolor}
\usepackage{dsfont}
\usepackage{booktabs}
\usepackage[labelfont=bf]{caption}
\usepackage{wrapfig}
\usepackage{algpseudocode}
\usepackage{tocloft}
\usepackage[nottoc]{tocbibind} %Bibliography in ToC

\let\OLDthebibliography\thebibliography %% bibliography printout settings
\renewcommand\thebibliography[1]{
  \OLDthebibliography{#1}
  \setlength{\itemsep}{.4ex}
}


\usepackage[capitalise,nameinlink,noabbrev]{cleveref}

\usepackage{tikz}
\usetikzlibrary{positioning,fit}
\usepackage{tikz-3dplot}

%=== Theorem environments (registered with thmtools) ===
\declaretheorem[name=Theorem]{theorem}
\declaretheorem[name=Lemma]{lemma}
\declaretheorem[name=Proposition]{proposition}
\declaretheorem[name=Corollary]{corollary}

\declaretheorem[name=Conjecture]{conjecture}
\declaretheorem[name=Fact]{fact}
\declaretheorem[name=Assumption]{assumption}
\declaretheorem[name=Claim]{claim}

\theoremstyle{definition}
\declaretheorem[name=Definition]{definition}
\declaretheorem[name=Problem]{question}
\declaretheorem[name=Example]{example}

\theoremstyle{remark}
\declaretheorem[name=Remark]{remark}


\usepackage{thm-restate}

% Bold label in example environments
\makeatletter
\def\th@example{%
  \thm@notefont{}% same as heading font
  \normalfont % body font
}
\def\th@definition{%
  \thm@notefont{}% same as heading font
  \normalfont % body font
}
\makeatother


% =============================================

\DeclareMathOperator{\rk}{rk}
\DeclareMathOperator{\crk}{\overline{rk}}
\DeclareMathOperator{\sk}{sk}
\DeclareMathOperator{\csk}{\overline{sk}}
\DeclareMathOperator{\mar}{mar}
\DeclareMathOperator{\Rcc}{R}
\DeclareMathOperator{\Ucc}{U}
\DeclareMathOperator{\Dcc}{D}
\DeclareMathOperator{\Ccc}{C}
\DeclareMathOperator{\dist}{dist}



\DeclareMathOperator{\res}{res}
\DeclareMathOperator{\supp}{supp}

\renewcommand{\P}{\operatorname{P}}


\DeclareMathOperator{\EQ}{\textsc{Eq}}
\DeclareMathOperator{\HD}{\textsc{HD}}
\let\SS\undefined
\DeclareMathOperator{\SS}{\textsc{SS}}

\DeclareMathOperator{\xor}{\textsc{Xor}}
\DeclareMathOperator{\Rank}{\textsc{Rank}}
\DeclareMathOperator{\HW}{\textsc{HW}}
\DeclareMathOperator{\Parity}{\textsc{Parity}}
\DeclareMathOperator{\PRank}{\textsc{PRank}}

\DeclareMathOperator{\SK}{\textsc{SK}}

\newtheorem{step}{Step}


\crefname{example}{Example}{Examples}


\newcommand{\rankOne}{\Rank_{\leq 1}}

\title{\vspace{-3cm}Rank-One\vspace{-1cm}}

% \setlength{\parindent}{0pt}

\newcommand{\valentin}[1]{{\color{red}\textbf{VI:} #1 }}
\newcommand{\mika}[1]{{\color{Cerulean}\textbf{Mika:} #1 }}
\newcommand{\kaave}[1]{{\color{blue}\textbf{Kaave:} #1 }}
\newcommand{\anastasia}[1]{{\color{DarkOrchid}\textbf{Anastasia:} #1 }}

\usepackage[most]{tcolorbox}
\newtcbox{\dominobox}{
  on line,
  boxrule=0.5pt,
  arc=3pt,
  left=3pt,right=3pt,top=0pt,bottom=0pt,
  boxsep=0pt
}

\newcommand{\domino}[2]{%
  \,\dominobox{$\begin{matrix} #1 \\ #2 \end{matrix}$}\,%
}

%STANDARD COMMANDS
\newcommand{\N}{\mathbb{N}}
\newcommand{\Z}{\mathbb{Z}}
\newcommand{\R}{\mathbb{R}}
\newcommand{\C}{\mathbb{C}}
\newcommand{\Q}{\mathbb{Q}}
\newcommand{\T}{\mathbb{T}}
\newcommand{\One}{\mathds{1}}
\newcommand{\cP}{\mathcal{P}}
\newcommand{\cL}{\mathcal{L}}
\newcommand{\cI}{\mathcal{I}}
\newcommand{\cD}{\mathcal{D}}
\newcommand{\cX}{\mathcal{X}}
\newcommand{\cY}{\mathcal{Y}}
\newcommand{\cZ}{\mathcal{Z}}


%RENEWED COMMANDS
\renewcommand{\epsilon}{\varepsilon}
\renewcommand{\leq}{\leqslant}
\renewcommand{\geq}{\geqslant}
\renewcommand{\setminus}{\backslash}
\renewcommand{\Re}{{\rm Re}}
\renewcommand{\Im}{{\rm Im}}
\renewcommand{\subset}{\subseteq}
\renewcommand{\supset}{\supseteq}

\usetikzlibrary{
    babel,
    arrows,
	calc,
	patterns,
	intersections,
    backgrounds,
    trees,
    mindmap,
    fadings,
	decorations.pathreplacing,
	decorations.pathmorphing,
	decorations.text,
	decorations.markings,
    scopes,
	shapes,
	positioning,
    arrows.meta,
    math,
    shapes.geometric,
    calligraphy
}


% COLOURED BOXES


%\definecolor{resultcolor}{HTML}{FFE1D9}
\definecolor{resultcolor}{HTML}{EED7FC}
\definecolor{definitioncolor}{HTML}{FCDFBE}
%\definecolor{definitioncolor}{HTML}{EFF9BD}
%\definecolor{definitioncolor}{HTML}{FFE4B8}
\definecolor{questioncolor}{HTML}{DDE8EB}

\newenvironment{boxtheorem}{\begin{theorem}}{\end{theorem}}
\tcolorboxenvironment{boxtheorem}{colback=resultcolor, colframe=white, colbacktitle=resultcolor, coltitle=resultcolor}

\newenvironment{boxlemma}{\begin{lemma}}{\end{lemma}}
\tcolorboxenvironment{boxlemma}{colback=resultcolor, colframe=white, colbacktitle=resultcolor, coltitle=resultcolor}

\newenvironment{boxproposition}{\begin{proposition}}{\end{proposition}}
\tcolorboxenvironment{boxproposition}{colback=resultcolor, colframe=white, colbacktitle=resultcolor, coltitle=resultcolor}

\newenvironment{boxcorollary}{\begin{corollary}}{\end{corollary}}
\tcolorboxenvironment{boxcorollary}{colback=resultcolor, colframe=white, colbacktitle=resultcolor, coltitle=resultcolor}

\newenvironment{boxdefinition}{\begin{definition}}{\end{definition}}
\tcolorboxenvironment{boxdefinition}{colback=definitioncolor, colframe=white, colbacktitle=definitioncolor, coltitle=definitioncolor}

\newenvironment{boxquestion}{\begin{question}}{\end{question}}
\tcolorboxenvironment{boxquestion}{colback=questioncolor, colframe=white,
    colbacktitle=questioncolor, coltitle=questioncolor}


 